\section{Conclusion}
\label{sec:conclusion}

QK-LUTFormer identifies two conditions for lookup-addressable spiking
attention. First, the internal discrete state must preserve semantic
address--response alignment: correct Q/K tuples predict held-out responses,
whereas coarse and shuffled controls do not recover the same structure.
Second, the retrieved response must preserve the dynamic current contract seen
by downstream BN/LIF state. Global-residual backoff makes unsupported lookup
explicit, subspace-decoupled residual tables retain most full-address gain
under a bounded entry budget, and TCSLU restores time/channel statistics for
end-to-end current replacement.

The completed QKFormer experiments support this chain across datasets, time
steps, calibration subsets, and independent checkpoints. The QK-contract
Spikformer intervention further shows that lookup fidelity depends on the
attention interaction, while unmodified Spikformer defines a useful mismatch
boundary. On CIFAR10-DVS, the best LUT-only deployment reaches $83.9\%$, and
a SpiLiFormer transfer reaches $81.7\%$ from its paired $81.2\%$ teacher with
the original projection bypassed. Because the from-scratch SpiLiFormer did not
reproduce the published clean accuracy and neither event route passed the
$84.0\%$ gate, this evidence remains bounded. Overall, the findings support a
compact and auditable framework for QK-compatible spiking attention. They do
not establish a universal principle for spiking neurons, measured hardware
efficiency, ImageNet performance, strong event-data accuracy, unrestricted
architecture transfer, or a complete production full-graph replacement.
